\documentclass[12pt]{article}
\usepackage{amsmath,amssymb,amsthm,hyperref}

\newtheorem{theorem}{Theorem}
\newtheorem{corollary}[theorem]{Corollary}
\newtheorem{remark}[theorem]{Remark}
\newtheorem{definition}[theorem]{Definition}

\title{Distances Occurring Between One and $n$ Times}
\author{}
\date{May 2026}

\begin{document}
\maketitle

\begin{abstract}
Let $A \subset \mathbb{R}^2$ be a set of $n$ points. We study distances determined by $A$ whose multiplicity lies between $1$ and $n$. We give a complete unconditional disproof of the literal statement that there must always be two such distances: an equilateral triangle determines only one distance. We also record the actual status of the stronger asymptotic question in the literature: known results on large distances do not yield the bound needed for a proof that $g(n) \to \infty$.
\end{abstract}

\section{Statement of the problem}

For a finite set $A \subset \mathbb{R}^2$, let
\[
\Delta(A) := \{ |x-y| : x,y \in A,\ x \neq y \}
\]
be the set of distinct distances determined by $A$. For $d \in \Delta(A)$, write $m_A(d)$ for the number of unordered pairs $\{x,y\} \subset A$ with $|x-y| = d$.

We are interested in distances $d \in \Delta(A)$ satisfying
\[
1 \le m_A(d) \le n.
\]

Define
\[
L(A) := \#\{ d \in \Delta(A) : 1 \le m_A(d) \le n \}.
\]
For each $n$, define the extremal quantity
\[
g(n) := \min\{ L(A) : A \subset \mathbb{R}^2,\ |A|=n \}.
\]

The original questions are:
\begin{enumerate}
\item Must $g(n) \ge 2$?
\item Must $g(n) \to \infty$ as $n \to \infty$?
\end{enumerate}

\section{A small counterexample}

\begin{theorem}
The assertion $g(n) \ge 2$ is false in general.
\end{theorem}

\begin{proof}
Take $n=3$ and let $A$ be the vertex set of an equilateral triangle. Then $\Delta(A)$ has exactly one element, namely the side length of the triangle, and that distance occurs exactly three times. Hence $L(A)=1$, so $g(3)=1<2$.
\end{proof}

Thus the first question has a negative answer if interpreted literally for every $n$.

\section{What the literature does and does not give}

The argument in an earlier draft appealed to the following strong assertion: for each fixed $k$, the $k$ largest distinct distances determined by a planar $n$-point set each occur at most $n$ times. If this were available, then the asymptotic conclusions would follow immediately. However, this is not what the classical papers prove.

For clarity, define the following assertion.

\begin{quote}
For every fixed integer $k \ge 1$, there exists $N(k)$ such that for every $n \ge N(k)$ and every set $A \subset \mathbb{R}^2$ with $|A|=n$, if
\[
\delta_1 > \delta_2 > \cdots > \delta_s
\]
are the distinct distances determined by $A$ in decreasing order, then
\[
m_A(\delta_i) \le n \qquad (1 \le i \le k).
\]
\end{quote}

If this assertion is accepted, then for each fixed $k$ the $k$ largest distances all contribute to $L(A)$, and hence $g(n) \ge k$ for all sufficiently large $n$. In particular one would obtain both eventual validity of $g(n) \ge 2$ and the divergence $g(n) \to \infty$.

What is actually available online from standard references is weaker.

\begin{enumerate}
\item In \emph{On large distances in planar sets} (Discrete Mathematics, 1987), Vesztergombi studied the multiplicity of the second-largest distance among $n$ points in the plane.
Later survey and expository sources summarizing this paper state that the second-largest distance can occur at most $\tfrac{3}{2}n$ times.

\item In later summaries of the large-distance literature, one finds the statement that the number of occurrences of the $k$-th largest distance in the plane is always smaller than $2kn$.
This is again far weaker than the bound $n$ required above.
\end{enumerate}

Thus the statement needed by the earlier draft is not supplied by these references. In particular, the known bound $m_A(\delta_2) \le \tfrac{3}{2}n$ does not even imply that the second-largest distance contributes to $L(A)$, because our threshold is exactly $n$.

Therefore the asymptotic claims are not merely missing a citation in the present note: they require genuinely stronger input than the classical bounds quoted above.

\begin{remark}
There is therefore a sharp distinction between the two parts of the original problem.
The first question has a complete unconditional negative answer, proved above.
The stronger asymptotic claims require additional input not currently justified by the standard large-distance references discussed here.
\end{remark}

\section{Conclusion}

The present note proves the following unconditional statement.

\begin{theorem}
Let $g(n)$ be defined as above.
\begin{enumerate}
\item The inequality $g(n) \ge 2$ is false for all $n$, since $g(3)=1$.
\item No unconditional proof of the asymptotic statements is given in this note, and the standard large-distance bounds quoted in the literature do not suffice to prove them.
\end{enumerate}
\end{theorem}

\begin{thebibliography}{9}

\bibitem{V87}
K.~Vesztergombi,
\emph{On large distances in planar sets},
Discrete Mathematics \textbf{67} (1987), 191--198.

\bibitem{ELV89}
P.~Erd\H{o}s, L.~Lov\'asz, and K.~Vesztergombi,
\emph{On the graph of large distances},
Discrete \& Computational Geometry \textbf{4} (1989), 541--549.

\bibitem{MP13}
F.~Mori\'c and J.~Pach,
\emph{Large simplices determined by finite point sets},
Beitr\"age zur Algebra und Geometrie \textbf{54} (2013), 481--487.

\end{thebibliography}

\end{document}